%% kapa: exam=1 round=24 year=2013 period=1 form=A subject=civil_law subject_ko=민법 seq=31 number=31
%% source: Q-Net 감정평가사 기출문제 (cst00309) artlSeq=5073802 file=extracted/1008/5073802/2013년 제24회 감정평가사 1차시험 기출문제/1교시(A형).hwp
%% conversion: kordoc 4.24.51 -> tools/gen1.py (auto); review=proofread-sonnet; flags=none
\begin{question}{31}
지상권에 관한 설명으로 옳지 않은 것은? (다툼이 있으면 판례에 의함)
\begin{choices}
  \item 지료는 지상권의 성립요소가 아니다.
  \item 토지에 저당권을 취득한 자가 그 토지 위에 차후 용익권이 설정되어 담보가치가 떨어지는 것을 막기 위하여 자신의 명의로 지상권을 설정한 경우 저당권의 피담보채권이 변제로 소멸하면 지상권도 소멸한다.
  \item 지상권자가 건물을 축조한 뒤 지상권을 유보한 채 건물만을 양도할 수도 있고 건물 소유권을 유보한 채 지상권만을 양도할 수도 있다.
  \item 기존 건물을 소유하기 위해 지상권을 설정받은 경우 지상권의 최단존속기간에 관한 민법규정이 적용된다.
  \item 존속기간을 영구(永久)로 약정한 지상권도 유효하다.
\end{choices}
\end{question}
\answer{4}
